\documentclass[11pt]{article}

\usepackage[margin=1in]{geometry}
\usepackage{booktabs}
\usepackage{siunitx}
\usepackage{amsmath}

\title{Assignment 1: Branch Predictor Implementation}
\author{Course 6385}
\date{\today}

\begin{document}
\maketitle

\section{How to run}

From \texttt{src}, build the driver and score one trace. The program prints
one line, the mispredictions per thousand instructions for that file.

\begin{verbatim}
cd src
make
./predict ../traces/164.gzip/gzip.trace.bz2
\end{verbatim}

Score every trace from the repository root. \texttt{run} calls
\texttt{src/predict} on each file under \texttt{traces}. Build first.
The trace reader decompresses each file with \texttt{/usr/bin/bzip2 -dc}.

\begin{verbatim}
csh ./run traces
\end{verbatim}

Check the Appendix~B worksheet from \texttt{src}. A match ends with
\texttt{PASS: 22 / 22 rows match the csv tables}.

\begin{verbatim}
cd src
make appendix_b_check
./appendix_b_check
\end{verbatim}

\let\oldsection\section
\renewcommand{\section}{\clearpage\oldsection}

\section{What was implemented}

The submitted predictor guesses whether a conditional branch is taken.
\texttt{my\_predictor.h} has two copies of the same idea.
\texttt{pm\_predictor} is the tiny Appendix~B example: 16 local counters,
4 global sets, 2 ways, and 4 bits of history. The checker uses only that
copy. \texttt{scaled\_pm\_predictor} is the copy \texttt{predict.cc} runs on
the traces. It is the Pentium~M size: a local table, a history of recent
outcomes, and a tagged global table with 4 ways.

For each conditional branch the large predictor does three things.

\begin{enumerate}
  \item \textbf{Local guess.} A table of 4096 two-bit counters is indexed by
        \texttt{PC[11:0]}. Each counter counts how often that address was
        taken. Values 0 and 1 predict not taken. Values 2 and 3 predict
        taken. The counter stops at 0 and at 3.
  \item \textbf{Recent outcomes.} A 15-bit register shifts in a 1 when the
        branch was taken and a 0 when it was not. This replaces the Pentium~M
        path register, which the assignment does not require.
  \item \textbf{Global guess.} \texttt{PC[16:2]} is XORed with those 15
        history bits, the same split as Appendix~B with a wider address.
        Bits \texttt{[14:6]} pick one of 512 sets. Bits \texttt{[5:0]} are a
        tag. Each set holds four entries. If one entry is valid and the tag
        matches, its counter makes the prediction. Otherwise the local guess
        is used.
\end{enumerate}

The global table is updated only when it is useful. The local counter is
trained when the global table misses. A new global entry is created only
when that local guess was wrong, in an empty slot or in place of the
least-recently-used one. A replaced counter starts at 2 if the branch was
taken and at 1 if it was not. If a global hit is wrong and the local guess
was right, that entry is deleted.

The trace tables start at zero. The worksheet copy keeps the Appendix~B
starting values, so its counters still match the example.

\subsection{What was skipped}

The loop predictor, the branch target buffer, and the path register are
omitted. Unconditional branches are predicted taken and do not change the
tables. Branch targets are not predicted. The score is direction misses only.

\section{Check against Appendix~B}

Appendix~B walks through 22 branches at \texttt{0x44}, \texttt{0x6C}, and
\texttt{0x90}. The expected history, index, local counters, and global
entries are in \texttt{hybrid.csv}, \texttt{local.csv}, and
\texttt{global.csv}. \texttt{appendix\_b\_check} feeds each address and
outcome to \texttt{pm\_predictor} and compares every saved field.

That run matches 22 of 22 rows. The large predictor is a different width, so
it is not checked against those rows.

\section{Trace results}

Each of the 20 traces is 100 million instructions. The score is
mispredictions per thousand instructions (MPKI): 6.311 means about six
wrong direction guesses in every thousand instructions. ``Always taken''
guesses that every conditional branch is taken. ``gshare'' is the
32,768-entry predictor shipped with the simulator. ``Hybrid'' is
\texttt{scaled\_pm\_predictor}.

\begin{table}[ht]
\centering
\begin{tabular}{l S[table-format=3.3] S[table-format=2.3] S[table-format=2.3]}
\toprule
{Trace} & {Always taken} & {gshare} & {Hybrid} \\
\midrule
164.gzip      & 70.025 & 12.473 & 14.543 \\
175.vpr       & 85.312 & 13.415 & 13.693 \\
176.gcc       & 86.071 & 11.254 & 9.126 \\
181.mcf       & 102.497 & 15.837 & 16.920 \\
186.crafty    & 48.651 & 5.837 & 5.129 \\
197.parser    & 78.421 & 10.008 & 8.995 \\
201.compress  & 76.983 & 7.831 & 8.978 \\
202.jess      & 92.343 & 1.562 & 1.148 \\
205.raytrace  & 109.506 & 2.756 & 2.139 \\
209.db        & 91.451 & 3.909 & 3.824 \\
213.javac     & 109.644 & 2.267 & 2.382 \\
222.mpegaudio & 110.339 & 2.188 & 2.400 \\
227.mtrt      & 112.022 & 2.657 & 2.314 \\
228.jack      & 77.958 & 3.033 & 2.538 \\
252.eon       & 24.226 & 1.807 & 1.730 \\
253.perlbmk   & 66.833 & 2.554 & 1.905 \\
254.gap       & 75.950 & 3.926 & 4.554 \\
255.vortex    & 70.524 & 1.222 & 0.854 \\
256.bzip2     & 178.257 & 0.094 & 0.071 \\
300.twolf     & 59.221 & 21.489 & 22.973 \\
\midrule
{Average}     & 86.311 & 6.306 & 6.311 \\
\bottomrule
\end{tabular}
\caption{Wrong direction guesses per thousand instructions.}
\end{table}

Always taken averages 86.311. The hybrid averages 6.311, and gshare averages
6.306. The two accurate predictors differ by 0.005. gcc and perlbmk are
better on the hybrid. gzip and twolf are better on gshare.

The accuracy matches because the global table is reserved for branches the
local table gets wrong. A stored global guess therefore does not replace a
local guess that was already right.

\begin{table}[ht]
\centering
\begin{tabular}{l r r}
\toprule
& gshare & Hybrid \\
\midrule
Stored bits  & 65,551 & 30,735 \\
Physical rows & 256 & 64 and 64 \\
Average MPKI & 6.306  & 6.311 \\
\bottomrule
\end{tabular}
\caption{Size of the two predictors. Physical rows are after the bits are folded into a near-square array. One row is read.}
\end{table}

gshare stores 32,768 two-bit counters and a 15-bit history, 65,551 bits.
The hybrid stores 4,096 local counters (8,192 bits), 2,048 global entries
with a 6-bit tag, a 2-bit counter, and a valid bit (18,432 bits), a 2-bit
age on each entry (4,096 bits), and the same 15-bit history. That is
30,735 bits. gshare holds 2.13 times as much state.

A wrong guess costs the same in either design, because the miss rates match.
Area and idle power follow the stored bits, so both are 2.13 times lower
for the hybrid.

Lookup delay is estimated for a Pentium~M Banias, the 130\,nm part at
1.3\,GHz. One cycle is 769\,ps. The fanout-of-4 inverter delay used here is
25.2\,ps, so a cycle is about 30 FO4.\footnote{Banias clock and FO4 delay
from the Stanford CPU DB. The gate delays below are taken from a 130\,nm
measurement of a 64-row memory: a 6-bit decoder plus wordline is 5 FO4, and
the bitline plus sense amplifier is 2.4 FO4.}
The memories are not one entry per row. The same bits are folded so the
array is near a square, and the address bits that no longer pick a row
become a column multiplexer. gshare's 64\,Kbit array is 256 rows by 256
columns, with a 128-to-1 column mux. The local table is 64 rows by 128
columns, with a 64-to-1 mux. Each global way is its own 64-row memory, and
the four ways are read together.

The terms, in FO4 delays, are a parallel XOR \(T_{\oplus}=2\), a decoder
\(T_{\mathrm{dec}}(n)=5+(n-6)\) for an \(n\)-bit row address, a bitline and
sense amplifier \(T_{\mathrm{bit}}(R)=2.4\times R/64\), a column mux
\(T_{\mathrm{col}}(k)=\frac{1}{2}\log_2 k\), four 6-bit tag comparisons
\(T_{\mathrm{cmp}}=3\), a way mux \(T_{\mathrm{way}}=2\), and the final
local-or-global mux \(T_{\mathrm{sel}}=2\). One row is driven.
\(T_{\mathrm{bit}}\) grows with the number of rows because the idle rows
still load the bitline.
\begin{align}
L_{\mathrm{gshare}} &= T_{\oplus}+T_{\mathrm{dec}}(8)+T_{\mathrm{bit}}(256)+T_{\mathrm{col}}(128) \nonumber\\
&= 2+7+9.6+3.5 = 22.1~\text{FO4}, \\
L_{\mathrm{local}} &= T_{\mathrm{dec}}(6)+T_{\mathrm{bit}}(64)+T_{\mathrm{col}}(64) \nonumber\\
&= 5+2.4+3 = 10.4~\text{FO4}, \\
L_{\mathrm{global}} &= T_{\oplus}+T_{\mathrm{dec}}(6)+T_{\mathrm{bit}}(64)+T_{\mathrm{col}}(8)+T_{\mathrm{cmp}}+T_{\mathrm{way}} \nonumber\\
&= 2+5+2.4+1.5+3+2 = 15.9~\text{FO4}, \\
L_{\mathrm{hybrid}} &= \max(L_{\mathrm{local}},L_{\mathrm{global}})+T_{\mathrm{sel}} \nonumber\\
&= 15.9+2 = 17.9~\text{FO4}.
\end{align}
At 25.2\,ps per FO4 this is 557\,ps for gshare and 451\,ps for the hybrid.
Both fit in the 769\,ps cycle. The hybrid lookup is 4.2 FO4, about 106\,ps,
shorter. The difference is the bitline: 9.6 FO4 at 256 rows against 2.4 FO4
at 64 rows. The tag compare and the two muxes add 7 FO4, which does not use
up that whole saving. This is a sum of those assumptions, not a measured
silicon delay.

\section{Conclusion}

The small predictor matches the Appendix~B example on all 22 rows. The
trace predictor is the Pentium~M hybrid at full size: 4,096 local counters,
15 bits of outcome history, and 512 sets of four tagged counters. On the
traces it cuts the always-taken miss rate from 86.311 to 6.311~MPKI, level
with gshare at 6.306, in 2.13 times less storage. On the Banias estimate
the lookup is 17.9 FO4 against 22.1 FO4 for gshare, about 106\,ps inside the
same cycle.

\end{document}
